%=========== ΛΥΣΗ - 14ο ΘΕΜΑ Α ===========
\providecommand{\theoriaref}[3]{%
\begin{center}
\fcolorbox{maincolor!70!black}{maincolor!8}{%
\begin{minipage}{.88\linewidth}
\textcolor{maincolor!80!black}{\kerkissans{\textbf{#1~\ref{#2}}}}\\[1mm]
#3
\end{minipage}}
\end{center}}
\providecommand{\orismosref}[1]{%
\theoriaref{Ορισμός}{#1}{Η απάντηση δίνεται από τον αντίστοιχο ορισμό.}}
\providecommand{\apodeixiref}[1]{%
\theoriaref{Θεώρημα}{#1}{Η ζητούμενη απόδειξη δίνεται στην αντίστοιχη απόδειξη.}}

\begin{thema}{Α}
\begin{erwthma}
\item \apodeixiref{thewrhma:paragwgos_efaptomenhs}

\item \orismosref{orismos:efaptomenh_grafikhs_parastashs}

\item \orismosref{orismos:antistrofh_synarthsh}

\item Οι χαρακτηρισμοί των προτάσεων είναι:
\begin{alist}
\item \textbf{Σωστό}. Αν η $f$ είναι $1-1$ στο $[a,\beta]$, τότε δεν μπορεί να ισχύει $f(a)=f(\beta)$. Άρα δεν ικανοποιεί όλες τις προϋποθέσεις του θεωρήματος \eng{Rolle}.
\item \textbf{Σωστό}. Είναι ο τύπος ολοκλήρωσης κατά παράγοντες:
\[ \dint_a^\beta f'(x)g(x)\d x=[f(x)g(x)]_a^\beta-\dint_a^\beta f(x)g'(x)\d x. \]
\item \textbf{Σωστό}. Αν η συνεχής $f$ είναι μη αρνητική και το ολοκλήρωμά της στο $[a,\beta]$ είναι $0$, τότε δεν μπορεί να είναι θετική σε κανένα σημείο. Άρα $f(x)=0$ σε όλο το διάστημα.
\item \textbf{Σωστό}. Από το θεώρημα \eng{Fermat}, αν η $f$ έχει τοπικό ακρότατο σε εσωτερικό σημείο $x_0$ και είναι παραγωγίσιμη σε αυτό, τότε $f'(x_0)=0$.
\item \textbf{Λάθος}. Πολυωνυμική συνάρτηση βαθμού τουλάχιστον $2$ δεν έχει πλάγια ασύμπτωτη, αφού η διαφορά της από οποιαδήποτε ευθεία δεν τείνει στο $0$ στο άπειρο.
\end{alist}

\item Σωστή απάντηση είναι η \textbf{β}. Αν $f'(x_0)=0$ και $f''(x_0)>0$, τότε η $f$ παρουσιάζει τοπικό ελάχιστο στο $x_0$.
\end{erwthma}
\end{thema}
%=========== ΤΕΛΟΣ ΛΥΣΗΣ - 14ο ΘΕΜΑ Α ===========
